% =============================================================================
% Ablation Study
% =============================================================================

%----------------------------------------------------------------------
\section{Ablation Study}
\label{sec:ablation}
%----------------------------------------------------------------------

The comparative evaluation demonstrates that LU-STW delivers the strongest overall accuracy on Amazon, Yelp, and MovieLens-100K among the methods considered, as reflected by the lowest MAE in Table~8. A natural follow-up question is whether this advantage arises from the proposed cross-domain mechanism or from favorable hyperparameter choices. To address this question, we perform a component-wise ablation in which individual parts of LU-STW are removed while all remaining training settings are kept unchanged. All runs use five-fold cross-validation on the same domain pairs, with an identical matrix factorization configuration for every variant. The ablated settings are Target Only, w/o Hungarian Alignment, w/o Confidence Weighting, and w/o Fine-Tuning. Table~\ref{tab:ablation_mae} reports the results for each dataset; the full LU-STW entry in that table is the reference baseline, and $\Delta$MAE denotes the MAE increase when a component is removed, where larger positive values indicate a more important contribution.

\begin{table}[htbp]
    \centering
    \footnotesize
    \setlength{\tabcolsep}{4pt}
    \caption{Ablation study on Amazon, Yelp, and MovieLens-100K. MAE: mean$\,\pm$\,std.\ over five folds; $\Delta$MAE relative to full LU-STW on the same dataset.}
    \label{tab:ablation_mae}
    \begin{tabular}{@{} l l c c @{}}
        \toprule
        Dataset & Variant & MAE & \makecell{$\Delta$MAE} \\
        \midrule
        \multirow{5}{*}{\makecell{Amazon}}
            & Target Only       & 0.7224$\,\pm$\,0.0103 & +0.0073 \\
            & w/o Hungarian     & 0.7210$\,\pm$\,0.0089 & +0.0059 \\
            & w/o Confidence    & 0.7208$\,\pm$\,0.0063 & +0.0057 \\
            & w/o Fine-Tuning   & 0.7202$\,\pm$\,0.0105 & +0.0051 \\
            & LU-STW            & 0.7151$\,\pm$\,0.0118 & 0.0000 \\
        \midrule
        \multirow{5}{*}{Yelp}
            & Target Only       & 0.8109$\,\pm$\,0.0325 & +0.0718 \\
            & w/o Hungarian     & 0.8060$\,\pm$\,0.0302 & +0.0669 \\
            & w/o Confidence    & 0.8104$\,\pm$\,0.0325 & +0.0713 \\
            & w/o Fine-Tuning   & 0.8109$\,\pm$\,0.0325 & +0.0718 \\
            & LU-STW            & 0.7391$\,\pm$\,0.0167 & 0.0000 \\
        \midrule
        \multirow{5}{*}{\makecell{MovieLens\\100K}}
            & Target Only       & 0.7431$\,\pm$\,0.0055 & +0.0385 \\
            & w/o Hungarian     & 0.7179$\,\pm$\,0.0054 & +0.0133 \\
            & w/o Confidence    & 0.7098$\,\pm$\,0.0052 & +0.0052 \\
            & w/o Fine-Tuning   & 0.7428$\,\pm$\,0.0056 & +0.0382 \\
            & LU-STW            & 0.7046$\,\pm$\,0.0017 & 0.0000 \\
        \bottomrule
    \end{tabular}
\end{table}

\subsection{Results Analysis}
\label{subsec:ablation_results}

The ablation results reveal a uniform trend: every ablated variant yields higher MAE and strictly positive $\Delta$MAE, whereas the complete LU-STW configuration (last row per dataset) recovers the reference performance. This confirms that the full model is Pareto-optimal and that no module is redundant.

The w/o Hungarian Alignment variant disables the Hungarian-based latent factor matching step in LU-STW; without this module, MAE increases on Amazon ($+0.0059$), Yelp ($+0.0669$), and MovieLens-100K ($+0.0133$) relative to the full model. The largest degradation on MovieLens reflects the need to align Romance and Drama latent factors across domains, while on Yelp the gap between $0.8060$ and $0.7391$ shows that even related categories (Veterinarians and Pet Services) benefit from optimal factor permutation via the Hungarian algorithm. These results confirm that Hungarian alignment reduces negative transfer and provides a sound initialization for subsequent fine-tuning.

The w/o Confidence Weighting variant removes user-level confidence scores during fine-tuning in LU-STW; this change raises MAE on Amazon ($+0.0057$), Yelp ($+0.0713$), and MovieLens-100K ($+0.0052$). On Yelp, where engagement patterns differ widely across users, applying the same update rate to every profile appears especially damaging, because users whose aligned and target representations agree should learn faster, while poorly matched users need gentler updates. On MovieLens, although the increase is smaller, the full model ($0.7046$) still clearly outperforms the ablated run ($0.7098 \pm 0.0052$), which supports confidence weighting as a way to transfer knowledge selectively rather than uniformly.

The w/o Fine-Tuning variant stops the target-domain adaptation stage after Hungarian alignment and blended initialization. Without fine-tuning, $\Delta$MAE reaches $+0.0382$ on MovieLens and $+0.0718$ on Yelp, showing that alignment by itself cannot replace personalization on the target domain. On Amazon, w/o Fine-Tuning ($0.7202$) remains above the full LU-STW result ($0.7151$), so the fine-tuning stage still adds measurable value even when source and target domains are relatively close.

The Target Only variant trains matrix factorization using target-domain ratings alone and does not use any source-domain knowledge, Hungarian alignment, confidence weighting, or fine-tuning from LU-STW. This baseline gives the largest $\Delta$MAE on Yelp ($+0.0718$, MAE $0.8109 \pm 0.0325$) and MovieLens-100K ($+0.0385$, MAE $0.7431 \pm 0.0055$), which shows that sparse target data alone cannot match the full cross-domain pipeline. On Amazon, Target Only reaches $0.7224 \pm 0.0103$ ($\Delta$MAE $+0.0073$), again worse than the complete model ($0.7151$). These outcomes indicate that LU-STW improves prediction by transferring and adapting source-domain structure, not by fitting the target domain in isolation.

Across all folds, the standard deviations remain small (for example, $0.0052$--$0.0056$ on MovieLens and $0.0063$--$0.0105$ on Amazon), so the ablation trends are stable and not tied to one lucky split. The size of each $\Delta$MAE varies by dataset---Yelp is most affected by removing confidence weighting or fine-tuning, MovieLens by removing Hungarian alignment or fine-tuning, and Amazon shows smaller but steady gaps---yet the same conclusion holds on every benchmark: each component of LU-STW contributes positively to the final recommendation accuracy.
